\subsection{The complete one- and two-object baseline}

\begin{theorem}[Proved: at most two objects require seventeen arrows]
\label{CP2:A-two-object-baseline-sixteen}
A finite category with at most two objects and classifying
space of type $\CP^2$ has at least seventeen arrows,
identities counted.
\end{theorem}
\begin{proof}
The one-object case follows from the complete minimum-monoid
classification, Theorem~\ref{A:monoid-minimum23-classification}.
A two-object model must be connected. If arrows exist in
both directions, Proposition~\ref{CP2:A-baseline-strong-two}
excludes even seventeen arrows. Otherwise it is a triangular
collage. By identity padding,
Lemma~\ref{CP2:A-identity-padding}, any model with fewer than
sixteen arrows would give one with exactly sixteen arrows,
with the same two objects and the same support. It therefore
suffices to exclude triangular collages at that exact total.

Write $(m,n,s)$ for the source corner order, target corner
order and cross-hom size. Thus $m+n+s=16$, with all three
entries positive. If both corners have zero rational second
cohomology, the relative-cup lemma excludes the projective-plane
square. Otherwise choose a corner with positive second
cohomology as the source, taking the opposite category if
necessary. The order-five corner classification gives
$m\geq5$. This choice does not assume $m\geq n$.

A singleton cross-hom realizes the join of the two connected
corner spaces and has zero rational second homology.
Two cross arrows are excluded by
Lemma~\ref{CP2:A-high-two-cross}; three by
Theorem~\ref{CP2:A-three-cross-reduction}. Hence $s\geq4$.

If $n=1$, the collage is an ordinary right act cone and
$m+s=15$. For $m\leq9$ use
Theorem~\ref{CP2:A-bounded-nine-acts}: the assumed
projective-plane type supplies bounded rational homology.
The only remaining possibilities are $(m,s)=(10,5)$ and
$(11,4)$. They are excluded, respectively, by
Proposition~\ref{CP2:A-ten-six-state-rational} and
Corollary~\ref{CP2:A-four-state-small}.

If $n=2$, both possible target monoids have projective
trivial right rational module. Indeed they are $C_2$ or
$\{1,z:z^2=z\}$; averaging in the first case and the
idempotent $z$ in the second give the splitting. The actual
orbit-act quotient functor is a rational homology and
cohomology-ring equivalence, by the exact-coinvariant
argument in Proposition~\ref{CP2:A-triangular-sixteen-finite-profiles}.
Thus it preserves the rational projective-plane target and
all-degree homological boundedness. Now $m+s=14$ and
$s\geq4$, so $m\leq10$. The bounded-nine theorem handles
$m\leq9$, and $(m,s)=(10,4)$ is covered by the
rational ten-element/six-state proposition. No actual
homotopy equivalence is inferred from this rational comparison.

It remains that $n\geq3$ and $s\geq4$. The possibilities
are exactly
\[
\begin{array}{c|l}
 m& (n,s)\\\hline
5&(3,8),(4,7),(5,6),(6,5),(7,4)\\
6&(3,7),(4,6),(5,5),(6,4)\\
7&(3,6),(4,5),(5,4)\\
8&(3,5),(4,4)\\
9&(3,4).
\end{array}
\]
The first row is Proposition~\ref{CP2:A-P5-source-sixteen},
since a positive-$H_2$ order-five source is $P_5$.
The next three rows are
Proposition~\ref{CP2:A-triangular-sixteen-finite-profiles},
including its separately proved graph-theorem treatment of
$(6,3,7)$. The last row is
Proposition~\ref{CP2:A-right-zero-four-states}, which includes
the projective opposite-corner alternatives explicitly.

For an arithmetic check independent of the table searches,
\path{monoid/code/two_object16_inventory.py} enumerates all
$55$ positive integer triples with $m\geq5$ and sum sixteen.
It verifies that the preceding gates are exhaustive and that
exactly the displayed fifteen profiles remain after the
small-cross and small-target reductions. The exact inventory
is \path{monoid/certificates/two_object16_inventory.json}.
Every case has now been excluded; padding gives the statement
for all smaller totals as well.
\end{proof}
